\section{Case studies: a controlled testbed and an external turbofan fleet}\label{sec:design}
\subsection{Why a controlled testbed}\label{sec:why}
Demonstrating an a-priori certificate honestly requires data in which the governing physics actually
controls degradation, so that the similitude departure $\boldsymbol\psi$ is meaningful and the
certificate can be checked rather than asserted. We therefore use a stochastic catalyst-deactivation
simulation~\cite{levenspiel1999chemical,perry2008handbook}, a process-engineering reliability problem
in which a catalytic unit ages as its catalyst deactivates and reaches end of life when activity
falls below a threshold; the prognostic task is to predict the remaining life for catalyst-replacement
scheduling. Dynamic similitude is \emph{genuine} in the simulator but is never assumed by the
estimator, which sees only noisy activity signals and a nominal rate model. This separation permits a
non-circular test: we impose a graded, physically meaningful departure from similitude and check that
the certificate degrades exactly as~\eqref{eq:cert} predicts.

\subsection{Simulator}\label{sec:sim}
Deactivation follows first order, $\mathrm{d}a/\mathrm{d}t=-k_{\mathrm{eff}}(T)\,a$, with failure at
activity $a\le0.3$. The primary Arrhenius channel is $k_1=A_1e^{-E_1/RT}$ ($E_1=80$~kJ/mol); an
optional secondary channel $k_2=A_2e^{-E_2/RT}$ of weight $\eta$ provides the controlled departure
($E_2\neq E_1$). Unit-to-unit variability enters as a lognormal pre-exponential, and the activity
signal carries additive noise. At $\eta=0$ the dimensionless clock $\tau=t\,k_1(T)$ collapses every
temperature onto $a(\tau)=e^{-\tau}$: exact similitude, with the characteristic life $T_L=1/k_1(T)$
spanning $10.2\times$ across the three temperatures (600/650/700~K). The weight $\eta$ thus moves the
testbed continuously from exact similitude ($\eta=0$) into progressively stronger violations.

\subsection{Predictor, scores, and reliability metrics}\label{sec:metrics}
The predictor extrapolates remaining life from the observed activity using the nominal rate, and the
conformal layer covers the residual. Each metric is read in maintenance terms. \emph{Coverage}
(target $1-\alpha=0.90$) is the fraction of intervals that contain the true RUL, i.e.\ the realised
reliability of the maintenance decision; the \emph{coverage gap} $\max(0,(1-\alpha)-\text{coverage})$
is the shortfall against the promised reliability. \emph{Efficiency} is the relative interval
back-off $b=q_T/\overline{\text{RUL}}_T$, which measures conservatism, that is the cost of
unnecessarily early replacement. With the one-sided score $V=\hat y(X)-Y$ the certified set is
$\{y\ge\hat y(X)-q_T\}$, so a planner acts on the certified fraction of predicted life $1-b$, and
we report it against three fixed bands: $b<0.5$ is
\emph{actionable}, more than half the predicted life being certified; $0.5\le b<1$ is
\emph{degraded}, a positive but shrinking lead time; and $b\ge1$ is \emph{degenerate}, the certified
lower bound being non-positive on average, so that the interval states only that failure has not
yet occurred and carries no replacement schedule. We also report the \emph{operational bound}
$2(a+L\lVert\Delta\boldsymbol\psi\rVert)$ and
the \emph{diagnostic} verdict. Calibration is \emph{unit-level}: every unit contributes one
score, taken at a single randomly chosen monitoring fraction, so the calibration set is
exchangeable across units. Numbers are averaged over five seeds, and coverage gaps are averaged
over (pair, seed) after the clip at zero, so an averaged gap is not the gap of the averaged
coverage; we keep the clipped convention throughout as the conservative reading. Three acceptance
criteria were fixed before running: (1) naive conformal miscovers while SCC recovers at $\eta=0$; (2) the a-priori
bound, with $L$ fit on a subset of (pair, $\eta$) configurations, upper-bounds the gap on
\emph{held-out} configurations; and (3) the coverage gap and the back-off degrade monotonically
with $\eta$. The three bands above are a reporting scale for the envelope, not a fourth criterion.

\subsection{Second case study: an externally authored turbofan fleet}\label{sec:cmapss}
The testbed above is ours, so its departure from similitude is a knob we set. To test the method
where the physics was authored elsewhere we add the NASA C-MAPSS turbofan
datasets~\cite{saxena2008cmapss,saxena2008damage}. FD002 contains 260 run-to-failure engines and
FD004 contains 249, and every engine visits all six operating regimes of the fleet, so each regime
carries roughly 260 units against the six bearings per condition available in the FEMTO benchmark,
a $43\times$ improvement in replication per condition.

The dimensionless coordinate is again derived rather than fitted. Gas-turbine performance collapses
onto referred (corrected) parameters built from the stagnation ratios $\theta=T_t/T_{\mathrm{std}}$
and $\delta_{\mathrm{p}}=P_t/P_{\mathrm{std}}$ (subscripted to keep it distinct from the scalar
departure $\delta$ of Section~\ref{sec:res-cert}), which follow from the recorded altitude, Mach
number and throttle setting through the standard atmosphere and the isentropic stagnation
relations; temperatures are referred by $\theta$, pressures by $\delta_{\mathrm{p}}$, speeds by
$\sqrt{\theta}$ and flows by $\delta_{\mathrm{p}}/\sqrt{\theta}$. The reduction is self-checking:
the sea-level
static regime returns $\theta=1.0000$ and $\delta_{\mathrm{p}}=0.9999$, and the implied damage clock
$\delta_{\mathrm{p}}/\sqrt{\theta}$ spans $3.5\times$ across the six regimes, so the regimes are far
apart in exactly the sense the method cares about.

Degradation is not the operating point but the departure from it, so the coordinate in which SCC
calibrates here is the deviation of each referred sensor from its regime's healthy baseline, the
standard construction of gas-path analysis. That baseline is taken over the first 20 cycles of
operation at the target regime, which a commissioned asset has, and uses no target failure data, so
the method's central claim is preserved. Engines are split disjointly into thirds for predictor
training, calibration and test, because every engine visits every regime and sharing engines would
leak dependence across the split; calibration takes one snapshot per engine per regime, matching the
unit-level scheme of Section~\ref{sec:metrics}. A ridge predictor is fitted on the source regime and
deployed on the target across all 30 ordered regime pairs at $\alpha=0.10$, with a random forest as
a capacity check.

Two departure quantities are computable before any target failure is observed, and both are tested
as candidates for $\boldsymbol\psi$: the distance between regimes in the referred ambient
coordinates $(\log\theta,\log\delta_{\mathrm{p}})$, and the distance between the regimes' healthy
gas-path signatures.
